\Needspace{8\baselineskip}\originalchapter{量子电动力学的可测过程}\label{ch:qed}
\originalsection{QED 顶点与图形}
QED 拉格朗日密度为 $\bar\psi(\ii\slashed D-m)\psi-\tfrac14F_{\mu\nu}F^{\mu\nu}$. 电子--光子顶点取 $-\ii e\gamma^\mu$, 电子传播子为 $\ii(\slashed q+m)/(q^2-m^2+\ii0)$, Feynman 规范中的光子传播子为 $-\ii g_{\mu\nu}/(q^2+\ii0)$. 外线光子附偏振矢量 $\epsilon_\mu$.

Compton 散射有两种顶点排列, 内部电子分别携带 $p+k$ 与 $p-k'$. $e^+e^-\to\gamma\gamma$ 也有两个电子交换排列. Møller 散射包含 $t,u$ 两通道且交换末态费米子带相对负号; Bhabha 散射包含 $s,t$ 两通道. 光子在 Abel 理论没有基本三光子或四光子顶点, light-by-light scattering 最低阶来自闭电子盒回路, 是四阶电荷振幅.

\originalsection{电子正电子湮灭}
在远高于 $m_e,m_\mu$ 而远低于需显著考虑 Z 交换的能区, $e^+e^-\to\mu^+\mu^-$ 的单光子振幅为
\begin{equation}
 \mathcal M=\frac{e^2}{s}[\bar v(p_2)\gamma^\mu u(p_1)]
 [\bar u(k_1)\gamma_\mu v(k_2)],
\end{equation}
整体符号依定义约定, 不影响单图概率. 两个初态自旋平均给 $1/4$, 用迹技巧得
\begin{equation}
 \overline{|\mathcal M|^2}=2e^4\frac{t^2+u^2}{s^2}
 =e^4(1+\cos^2\theta).
\end{equation}
最后一步利用无质量 $t=-s(1-\cos\theta)/2$, $u=-s(1+\cos\theta)/2$. 代相空间公式并用 $e^2=4\pi\alpha$,
\begin{equation}
 \frac{\dd\sigma}{\dd\Omega}=\frac{\alpha^2}{4s}(1+\cos^2\theta),\qquad
 \sigma=\frac{4\pi\alpha^2}{3s}.
\end{equation}
角分布反映中间自旋一的向量流, $1/s$ 标度来自无新质量尺度时的截面量纲. 阈值附近必须恢复 muon 质量, Z 极点附近还须加入弱中性流与干涉.

\originalsection{规范约束与可观测振幅}
外部光子偏振加上 $a k_\mu$ 不应改变物理振幅. 对电子线的两种 Compton 图, 将 $\epsilon\to k$ 并用外旋量方程运动, 两图中残留项相消. 单独删一张图通常破坏该抵消. 这说明规范不变性对完整图集合提出可复核条件, 不是画图之后的装饰性宣称.

真实横向光子只有两个物理偏振. 计算中用 $-g_{\mu\nu}$ 求和是借守恒电流消去非物理部分, 不是说有四个独立可测偏振. 自旋、规范和相空间必须在同一约定下组合.

\originalsection{真空极化与耦合定义}
带电虚粒子的真空极化修正改变光子两点函数, 使不同动量尺度探测到不同有效电荷. 一圈 QED 的带电费米子贡献导致高尺度有效 $\alpha$ 增大; 以常数质量与对数近似表示,
\begin{equation}
 \alpha(Q^2)\simeq\frac{\alpha(\mu^2)}{1-[\alpha(\mu^2)/(3\pi)]\log(Q^2/\mu^2)}
\end{equation}
只写一个单位电荷费米子且跨越阈值之外适用. 多物种贡献要乘电荷平方与简并因子. 裸参数和发散不是直接观测量, 必须指定实验或重整化条件定义有限参数. 本章解释其对截面的用途, 不重复完整 QFT 重整化体系.
